% Frozen source-aware proof packet, 2026-08-25.
%
% Authenticated primary-source dependency:
% Sebastien Boucksom and Mattias Jonsson, Global pluripotential theory over
% a trivially valued field, Ann. Fac. Sci. Toulouse 31 (2022), 647--836.
% Local primary PDF: tmp/pdfs/bj22_global_trivial.pdf.
% SHA-256: 8d5975337f4ac88a475929ea6a436ec4a15d08715d5582bc0ba0eb3e07152c52.
% Used exactly: Section 1.1.1, Definition 1.26, Proposition 1.28,
% Lemma 4.46, Section 11.3, Proposition 11.6(ii), Lemma 11.7, and
% Question 11.16.
%
% Scope firewall:
% * The center C is a linear projective line, even when dim X is larger.
%   Its dimension one is what makes the bad locus finite.
% * For n=2 the blowup along the Cartier divisor C has no exceptional
%   divisor, so v_C is explicitly defined to be ord_C.
% * The varying point valuations lie on pairwise different blowups, so no
%   one fixed dual complex contains a tail of the sequence.
% * No resolution of singularities is used.

\begin{proof}
All valuations in the statement are divisorial and hence have linear growth
by Proposition 1.28 of the cited source.

We first prove weak convergence.  Denote by \(\eta\) the generic point of
\(C\), and by \(\mathcal I\) the ideal sheaf of \(C\).  We use the following
elementary specialization fact.  If \(a\) is a nonzero regular function on
an open neighbourhood of \(\eta\), put
\[
 r_a=\operatorname{ord}_{\mathcal I_\eta}(a).
\]
For \(n\geq3\), the blowup construction identifies this number with
\(\operatorname{ord}_{E_C}(a)\); for \(n=2\), it is
\(\operatorname{ord}_C(a)\).  After shrinking around \(\eta\), the function
\(a\) lies in \(\mathcal I^{r_a}\), and its initial class
\[
 \operatorname{in}_C(a)
 \in \mathcal I^{r_a}/\mathcal I^{r_a+1}
 \simeq \operatorname{Sym}^{r_a}(\mathcal I/\mathcal I^2)
\tag{1}\label{eq:p11pn-initial-class}
\]
is a regular section that is nonzero at \(\eta\).  Its zero locus on the
curve \(C\) is therefore finite, after also adding the finite complement of
the chosen neighbourhood.

At every remaining point \(p\in C(k)\), the class in
\eqref{eq:p11pn-initial-class} has nonzero fibre.  Since \(C\) is a smooth
linear subspace of the smooth variety \(X\), the conormal injection
\[
 (\mathcal I/\mathcal I^2)\otimes k(p)
 \lhook\joinrel\longrightarrow
 \mathfrak m_p/\mathfrak m_p^2
\]
identifies the source with the transverse linear forms.  Its \(r_a\)-th
symmetric power is injective.  Thus the image of
\(\operatorname{in}_C(a)(p)\) in
\(\mathfrak m_p^{r_a}/\mathfrak m_p^{r_a+1}\) is nonzero.  Since
\(a\in\mathfrak m_p^{r_a}\), this proves
\[
             \operatorname{ord}_{\mathfrak m_p}(a)=r_a
\tag{2}\label{eq:p11pn-regular-specialization}
\]
outside a finite subset of \(C(k)\).

Now fix \(f\in k(X)^\times\).  On an affine open neighbourhood of \(\eta\),
write
\[
                              f=a/b
\]
with nonzero regular functions \(a,b\).  Apply
\eqref{eq:p11pn-regular-specialization} to both functions and discard the
union of their finite bad sets.  The exceptional valuation of the blowup at
a smooth point \(p\) is the \(\mathfrak m_p\)-adic order: the associated
graded ring is a polynomial ring over \(k(p)\), so this order is additive
and extends to the fraction field.  Hence, at every remaining point,
\[
 \begin{aligned}
  \operatorname{ord}_{E_p}(f)
    &=\operatorname{ord}_{\mathfrak m_p}(a)
      -\operatorname{ord}_{\mathfrak m_p}(b)\\
    &=r_a-r_b=v_C(f).
 \end{aligned}
\tag{3}\label{eq:p11pn-rational-specialization}
\]
Pairwise distinctness of the \(p_j\) makes the sequence eventually avoid
every finite subset of \(C(k)\).  Equation
\eqref{eq:p11pn-rational-specialization} therefore gives
\(v_j(f)=v_C(f)\) for all sufficiently large \(j\).  The weak topology on
the valuation locus is pointwise convergence on \(k(X)^\times\), so
\(v_j\to v_C\) weakly.

We next compute the exhaustion.  A nonzero homogeneous polynomial of degree
\(m\) has multiplicity at a point at most \(m\), and equality is attained by
the \(m\)-th power of a hyperplane through that point.  Definition 1.26 and
Lemma 4.46 of the cited source thus give
\[
                              T_\omega(v_j)=1.
\tag{4}\label{eq:p11pn-T-point}
\]
The homogeneous ideal of the linear line \(C\) is generated by linear
forms.  A nonzero degree-\(m\) homogeneous polynomial cannot lie in
\(\mathcal I_C^{m+1}\), while the \(m\)-th power of any hyperplane
containing \(C\) lies in \(\mathcal I_C^m\).  It follows, both for
\(\operatorname{ord}_{E_C}\) when \(n\geq3\) and for
\(\operatorname{ord}_C\) when \(n=2\), that
\[
                              T_\omega(v_C)=1.
\tag{5}\label{eq:p11pn-T-center}
\]

Finally, choose a hyperplane \(H_j\) through \(p_j\) that does not contain
\(C\), and let \(\ell_j\in H^0(X,L)\) be a defining section.  The hyperplane
has multiplicity one at \(p_j\), whereas \(\ell_j\) is nonzero at the
generic point of \(C\).  Therefore
\[
                          v_j(\ell_j)=1,
                          \qquad v_C(\ell_j)=0.
\tag{6}\label{eq:p11pn-moving-hyperplane}
\]
The section formula of Lemma 11.7 in the cited source gives
\[
 d_{\infty,\omega}(v_j,v_C)
 \geq \lvert v_j(\ell_j)-v_C(\ell_j)\rvert=1.
\tag{7}\label{eq:p11pn-strong-gap}
\]
Thus the sequence is not strongly convergent although it is weakly
convergent and its exhaustion values agree identically with the exhaustion
of its limit.  This is a negative answer to the implication asked for in
Question 11.16, in every dimension \(n\geq2\).
\end{proof}
